\documentclass[11pt,a4paper]{article}

% ---------- packages ----------
\usepackage[margin=25mm]{geometry}
\usepackage{amsmath}
\usepackage{booktabs}
\usepackage{graphicx}
\usepackage{siunitx}
\usepackage{xcolor}
\usepackage{caption}
\usepackage{listings}
\usepackage[hidelinks]{hyperref}

% ---------- helpers ----------
% Red placeholder text, delete each one once it is filled in
\newcommand{\todo}[1]{\textcolor{red}{[#1]}}

% Code listing style for the appendix
\lstset{
    basicstyle=\ttfamily\footnotesize,
    breaklines=true,
    frame=single,
    numbers=left,
    numberstyle=\tiny
}

% ---------- title ----------
\title{ENMT482 Manipulators Assignment\\[2mm]
       \large Automating Coffee Making with the UR5}
\author{
    Jaime Uy De Baron \\
    Josh Bergerson \\
    Johann Torres\\[2mm]
    Group 16, Cart 3
}
\date{\today}

\begin{document}
\maketitle

% ---------- abstract ----------
\begin{abstract}
The UR5 is a six-axis collaborative robot with a 5~kg payload capacity and 850~mm reach radius. This project aims to direct the UR5 to make a cup of coffee, from weighing out ground coffee to moving the finished coffee cup to a dedicated customer zone. 

This report will quantify the economic utility of the UR5 for this task by comparing it to a civilian performing the same procedure, considering cost and time efficiency alike. 
\end{abstract}

% =====================================================================
\section{Introduction}
% =====================================================================
\todo{State the aim. The UR5 performs the coffee making tasks (a to v) using
the interchangeable tools, and a time and motion study quantifies whether
automating the process is worthwhile.}

% =====================================================================
\section{Method}
% =====================================================================
% The brief does not say this section is required. Confirm with Chris or
% Rodney, and delete it if it is not expected.

\subsection{Locating the equipment}
\todo{Explain how the pose of each fixture relative to the robot base was found
from the points known in both frames. Include the planar case, the case with a
known tilt, and the case using two measured axes, and state which fixture used
which. Include one worked example with numbers, for instance the Rancilio
frame.}

\subsection{Tool and task poses}
\todo{Explain how the TCP pose was found for each task by chaining the fixture
pose, the task point in the fixture frame, and the tool offset. Show the chain
equation.}

\begin{equation}
{}^{UR}T_{TCP} = {}^{UR}T_{fixture} \; {}^{fixture}T_{task} \; \left({}^{TCP}T_{task}\right)^{-1}
\end{equation}

\subsection{Calibration and refinement}
\todo{List the values that were measured or tuned rather than calculated, and
why each was needed. Examples are the fixture fine tune offsets, approach angle
adjustments, lever pull angle, and jogged joint waypoints used to avoid
collisions.}

\begin{table}[h]
\centering
\caption{Values refined during testing.}
\begin{tabular}{llll}
\toprule
Task & Value & Setting & Reason \\
\midrule
\todo{b} & \todo{fixture fine tune} & \todo{(x, y, z) mm} & \todo{reason} \\
\todo{} & \todo{} & \todo{} & \todo{} \\
\bottomrule
\end{tabular}
\end{table}

\subsection{Program structure}
\todo{Briefly describe the file per task structure, the shared helper functions,
and how the tasks are chained in sequence.}

% =====================================================================
\section{Time and Motion Study}
% =====================================================================

\subsection{Method}
\todo{How the times were obtained. Human times come from the video of Rodney
making a coffee. Robot times come from running the finalised code. State how
each task was timed and any assumptions, such as whether tool changes are
included in the task time.}

\subsection{Results}

\begin{table}[h]
\centering
\caption{Human and robot time for each task.}
\begin{tabular}{clSS}
\toprule
Task & Description & {Human (\si{\second})} & {Robot (\si{\second})} \\
\midrule
a & Pick up the Rancilio tool and place it on the Mazzer Scale pan & 0 & 0 \\
b & Unlock the Mazzer Scale & 0 & 0 \\
c & Turn the Mazzer on, wait 15 s, turn it off & 0 & 0 \\
d & Pull the dosing lever until 20 g is deposited & 0 & 0 \\
e & Lock the Mazzer Scale & 0 & 0 \\
f & Remove the Rancilio tool from the Mazzer & 0 & 0 \\
g & Open the WDT fixture and place the Rancilio tool in & 0 & 0 \\
h & Release the Rancilio tool and close the WDT fixture & 0 & 0 \\
i & Turn the WDT rotor five revolutions & 0 & 0 \\
j & Open the WDT fixture, remove the tool, close the fixture & 0 & 0 \\
k & Place the Rancilio tool in the PUQ and wait 2 s & 0 & 0 \\
l & Remove from the PUQ and insert in the group head & 0 & 0 \\
m & Operate the cup dispenser & 0 & 0 \\
n & Pick up the cup and place it on the Rancilio Scale pan & 0 & 0 \\
o & Unlock the Rancilio Scale & 0 & 0 \\
p & Operate the hot water switch until 32 g is dispensed & 0 & 0 \\
q & Lock the Rancilio Scale & 0 & 0 \\
r & Place the cup of coffee in the customer zone & 0 & 0 \\
s & Remove the Rancilio tool from the group head & 0 & 0 \\
t & Actuate the silicone brush for 5 s & 0 & 0 \\
u & Actuate the bristle brush for 5 s & 0 & 0 \\
v & Return the Rancilio tool to the tool stand & 0 & 0 \\
\midrule
\multicolumn{2}{l}{\textbf{Total, complete process}} & 0 & 0 \\
\bottomrule
\end{tabular}
\end{table}

\todo{Comment on the results. Which tasks dominate the robot's time, and where
is the human faster or slower. Note tasks with fixed waits (15 s grind, 2 s
tamp, 5 s brushes) that neither can shorten.}

% =====================================================================
\section{Financial Analysis}
% =====================================================================

\subsection{Given figures}

\begin{table}[h]
\centering
\caption{Figures given in the assignment brief.}
\begin{tabular}{ll}
\toprule
Item & Value \\
\midrule
Cost to purchase the UR5 & \SI{40000}{\$} \\
Cost of tooling & \SI{5000}{\$} \\
Annual maintenance of the UR5 & \SI{400}{\$}, plus \SI{3}{\hour} shutdown \\
Hourly wage of an experienced barista & \SI{30}{\$\per\hour} \\
Breaks per \SI{8}{\hour} shift & 2 $\times$ \SI{15}{\minute} plus \SI{30}{\minute} lunch \\
Average price of a cup of coffee & \SI{5}{\$} \\
\bottomrule
\end{tabular}
\end{table}

\subsection{Assumptions}
% These are not given in the brief, so state your own and justify them
\todo{State the assumptions the analysis depends on, with a reason for each.
Examples are opening hours, operating days per year, demand in cups per hour
compared with each method's maximum output, and whether the robot needs a human
to refill cups or beans.}

\begin{table}[h]
\centering
\caption{Assumptions used in the analysis.}
\begin{tabular}{lll}
\toprule
Assumption & Value & Justification \\
\midrule
\todo{Operating hours per day} & \todo{} & \todo{} \\
\todo{Operating days per year} & \todo{} & \todo{} \\
\todo{Demand (cups per hour)} & \todo{} & \todo{} \\
\bottomrule
\end{tabular}
\end{table}

\subsection{Cost comparison}
\todo{Set out the calculation. Productive time per shift is the shift length
minus breaks, which is \SI{7}{\hour} of a \SI{8}{\hour} shift. Compute the
cups per hour each method can produce from the cycle times, the annual
cost of each, and the annual saving.}

\begin{align}
\text{Cups per hour} &= \frac{3600}{t_{\text{cycle}}} \\
\text{Annual saving} &= C_{\text{human, annual}} - C_{\text{robot, annual}} \\
\text{Payback period} &= \frac{C_{\text{UR5}} + C_{\text{tooling}}}{\text{Annual saving}}
\end{align}

\todo{Results, including the payback period, and a short table or figure.}

% =====================================================================
\section{Other Factors}
% =====================================================================
\todo{Discuss factors beyond the direct cost comparison.}

\begin{itemize}
    \item \textbf{Speed improvements:} \todo{for example faster joint moves,
    fewer tool changes, overlapping tasks, removing the fixed waits where
    possible}
    \item \textbf{Demand through the day:} \todo{what times people typically
    buy coffee and how that affects utilisation of the robot}
    \item \textbf{Other costs not included:} \todo{for example the coffee cart
    equipment, cup and bean refills, cleaning, supervision, downtime,
    insurance}
    \item \textbf{Quality and consistency:} \todo{dose and water accuracy of
    the robot compared with a human}
    \item \textbf{Safety and customer experience:} \todo{}
\end{itemize}

% =====================================================================
\section{Recommendation}
% =====================================================================
\todo{State clearly whether automating this task is worthwhile, and support it
with the payback period and the key factors above.}

% =====================================================================
\section{Conclusion}
% =====================================================================
\todo{Short summary of what was achieved and the recommendation.}

% ---------- references ----------
% \begin{thebibliography}{9}
% \bibitem{primer} ENMT482 Pose and Homogeneous Transforms primer.
% \bibitem{gitlab} UC-02-2024 Coffee Cart Modifications, GitLab.
% \end{thebibliography}

% ---------- appendix ----------
\appendix
\section{Code}
\todo{Include or link the task files if required. Delete this appendix if not.}

% \lstinputlisting[language=Python, caption={\texttt{\_partA.py}}]{_partA.py}

\end{document}
